\chapter{Menopause - managing symptoms without medication}

\section*{Definition and Overview}
Menopause is a natural biological transition marking the permanent cessation of menstruation, diagnosed retrospectively after 12 consecutive months of amenorrhoea. The decline in ovarian follicular activity leads to a dramatic reduction in circulating oestrogen levels, which can trigger a range of physical and psychological symptoms, including vasomotor symptoms (hot flushes and night sweats), sleep disturbances, mood changes, and vaginal dryness. While Menopausal Hormone Therapy (MHT) is the most effective pharmacological treatment, many women choose or are medically advised to avoid hormonal or non-hormonal medications. Managing menopausal symptoms without medication involves utilizing evidence-based lifestyle modifications, dietary adjustments, cognitive-behavioural strategies, and physical exercises to alleviate symptoms, support cardiovascular and bone health, and enhance overall quality of life.

\section*{Epidemiology}
Menopause typically occurs between the ages of 45 and 55 years, with the average age in developed countries, including Australia, being approximately 51 years. The transition phase (perimenopause) can begin several years prior and is characterized by irregular menstrual cycles. Vasomotor symptoms affect approximately 60\% to 80\% of women during the menopausal transition, varying in severity and duration. In approximately 20\% of women, symptoms are severe and can persist for seven years or longer. Due to concerns regarding the risks of MHT or personal preferences, a substantial proportion of women (up to 50\% to 70\% in some cohorts) seek non-pharmacological strategies to manage their symptoms, relying on self-care, complementary therapies, and lifestyle alterations.

\section*{Aetiology and Pathophysiology}
The symptoms of menopause are directly related to the depletion of ovarian follicles and the consequent loss of ovarian oestrogen and progesterone production. Non-pharmacological management aims to mitigate the physiological consequences of these endocrine changes through various biological and behavioural pathways:
\begin{itemize}
    \item \textbf{Vasomotor instability:} The loss of oestrogen destabilises the hypothalamic thermoregulatory zone, narrowing the sweating and shivering thresholds. Lifestyle modifications, such as layering clothing, using fans, and avoiding triggers (caffeine, alcohol, spicy foods), reduce the frequency of core body temperature spikes.
    \item \textbf{Cognitive-behavioural pathways:} Cognitive-Behavioural Therapy (CBT) helps restructure negative thoughts about symptoms, downregulating the autonomic nervous system arousal that exacerbates hot flushes and improves sleep hygiene.
    \item \textbf{Bone mineral density decline:} Accelerated bone resorption occurs due to oestrogen deficiency. Weight-bearing exercise and strength training stimulate osteoblast activity, helping to offset the risk of osteopenia and osteoporosis.
    \item \textbf{Cardiovascular risk increase:} Oestrogen has cardioprotective effects, and its decline adversely affects lipid profiles and vascular compliance. A Mediterranean-style diet high in dietary fibre and low in saturated fats helps maintain healthy lipid levels and endothelial function.
    \item \textbf{Genitourinary changes:} Urogenital atrophy occurs due to a lack of oestrogen, leading to vaginal dryness and dyspareunia. Non-hormonal vaginal lubricants and moisturisers restore moisture and maintain vaginal tissue pH, reducing friction during intercourse.
\end{itemize}

\section*{Clinical Features}
Women undergoing the menopausal transition exhibit a constellation of physical and psychological features that can be managed non-pharmacologically:
\begin{itemize}
    \item \textbf{Vasomotor symptoms:} Sudden sensations of heat, flushing of the face and chest, profuse sweating, and subsequent chills, which are often triggered by warm environments or stress.
    \item \textbf{Sleep disturbances:} Insomnia, frequent night waking (often secondary to night sweats), and daytime fatigue.
    \item \textbf{Mood swings and anxiety:} Irritability, low mood, anxiety, and difficulty concentrating (commonly referred to as "brain fog").
    \item \textbf{Genitourinary symptoms:} Vaginal dryness, burning, itching, dyspareunia, and urinary urgency.
    \item \textbf{Musculoskeletal symptoms:} Joint stiffness, muscle aches, and progressive loss of bone density.
    \item \textbf{Metabolic shifts:} Gradual weight gain, particularly abdominal fat deposition, and a slower metabolic rate.
\end{itemize}

\section*{Differential Diagnosis}
Before attributing symptoms entirely to menopause, other medical conditions must be ruled out:
\begin{itemize}
    \item \textbf{Thyroid disease:} Hyperthyroidism can mimic vasomotor symptoms, anxiety, and weight changes, while hypothyroidism can mimic fatigue, brain fog, and weight gain.
    \item \textbf{Clinical depression:} Persistent low mood and anhedonia that exceed typical menopausal mood swings, requiring targeted mental health intervention.
    \item \textbf{Sleep apnoea:} Primary sleep disorders presenting with daytime fatigue and night waking, which must be distinguished from menopause-related insomnia.
    \item \textbf{Rheumatological disorders:} Inflammatory joint diseases causing musculoskeletal pain and stiffness, which must be differentiated from menopausal arthralgia.
    \item \textbf{Cardiovascular disease:} Atypical chest pain or palpitations that should be thoroughly investigated rather than dismissed as menopausal anxiety.
\end{itemize}

\section*{When to Seek Medical Care}
Women should consult their doctor to discuss their menopausal transition and undergo routine health screenings.

\textbf{Call 000 (or local emergency services) for:}
\begin{itemize}
    \item \textbf{Severe chest pain or pressure:} Symptoms indicating acute coronary syndrome, as cardiovascular risk increases post-menopause.
    \item \textbf{Sudden shortness of breath:} Dyspnoea that may point to a pulmonary embolism or cardiac event.
    \item \textbf{Stroke symptoms:} Unilateral weakness, facial drooping, or speech difficulty, particularly in women with cardiovascular risk factors.
    \item \textbf{Severe psychiatric crisis:} Thoughts of self-harm or severe, debilitating depression.
    \item \textbf{Unexplained postmenopausal bleeding:} Any vaginal bleeding occurring 12 months or more after the final period must be investigated urgently to rule out endometrial cancer.
\end{itemize}

\section*{Diagnosis and Investigations}
Menopause is primarily a clinical diagnosis based on age and menstrual history. Investigations are directed toward ruling out other pathologies and establishing baseline health:
\begin{itemize}
    \item \textbf{Menstrual history review:} Documentation of irregular cycles followed by 12 months of amenorrhoea in a woman of appropriate age.
    \item \textbf{Thyroid-stimulating hormone (TSH):} To rule out hyperthyroidism or hypothyroidism as the cause of thermoregulatory, mood, or metabolic symptoms.
    \item \textbf{Lipid profile and fasting blood glucose:} To assess cardiovascular and metabolic risk, which increases during the menopausal transition.
    \item \textbf{Dual-energy X-ray absorptiometry (DEXA):} A bone density scan to screen for osteopenia or osteoporosis, particularly in women with risk factors.
    \item \textbf{Follicle-stimulating hormone (FSH) level:} Not routinely recommended for diagnosing menopause in women over 45 years, but may be measured in atypical or premature presentations (under 40 years).
\end{itemize}

\section*{Management}
Non-pharmacological management focuses on lifestyle adjustments, cognitive techniques, and physical therapies to alleviate symptoms and maintain long-term health:
\begin{itemize}
    \item \textbf{Cognitive Behavioural Therapy (CBT):} A structured programme focused on stress management, cognitive restructuring, and behavioural pacing, shown to reduce the impact and perceived severity of hot flushes and improve sleep.
    \item \textbf{Dietary modifications:} Adopting a Mediterranean diet rich in whole grains, vegetables, legumes, olive oil, and lean proteins; maintaining adequate calcium (1,300 mg daily for women over 50 years) and vitamin D intake to support bone health.
    \item \textbf{Physical exercise:} A balanced routine combining weight-bearing activities (such as brisk walking) and resistance training (such as weight lifting) at least 3 times weekly to maintain bone density and muscle mass, alongside aerobic exercise to support cardiovascular health.
    \item \textbf{Thermoregulation strategies:} Dressing in layers of natural fibres (cotton, silk), keeping indoor spaces cool with fans or air conditioning, and using cooling pillows or gel packs.
    \item \textbf{Sleep hygiene:} Establishing a regular sleep schedule, avoiding screens and heavy meals before bedtime, and keeping the bedroom dark and cool.
    \item \textbf{Non-hormonal vaginal moisturisers and lubricants:} Regular use of water- or silicone-based vaginal moisturisers to improve tissue hydration, and lubricants during intercourse to reduce discomfort.
    \item \textbf{Stress reduction techniques:} Mindfulness-based stress reduction (MBSR), yoga, or deep breathing exercises to lower stress hormones and reduce autonomic reactivity.
\end{itemize}

\section*{Complications}
Choosing a non-pharmacological approach is safe, but failure to address underlying health changes can lead to complications:
\begin{itemize}
    \item \textbf{Osteoporosis and fractures:} Severe bone loss leading to fragile bones and increased risk of minimal-trauma fractures (e.g. hip, wrist, spine).
    \item \textbf{Cardiovascular disease:} Elevated risk of coronary artery disease, stroke, and hypertension due to the loss of protective oestrogen.
    \item \textbf{Chronic insomnia and fatigue:} Persistent sleep disruption leading to cognitive impairment, mood disorders, and reduced work productivity.
    \item \textbf{Severe genitourinary syndrome of menopause (GSM):} Progressive vaginal atrophy, recurrent urinary tract infections, and complete avoidance of intimacy due to painful intercourse.
\end{itemize}

\section*{Prognosis and Outcomes}
The prognosis for managing menopausal symptoms without medication is highly positive for the majority of women. Lifestyle modifications and therapies like CBT have been shown to significantly reduce the impact of vasomotor symptoms and improve sleep quality. While these strategies do not restore oestrogen levels, they enhance the body's adaptation to the hormonal transition. Long-term adherence to exercise and dietary recommendations is highly effective in maintaining bone density and reducing cardiovascular risk, allowing women to transition through menopause with a high quality of life and minimal physical morbidity.

\section*{Prevention and Self-Care}
Self-care strategies focus on reducing symptom frequency and maintaining bone and cardiovascular health. Recommendations include:
\begin{itemize}
    \item \textbf{Avoiding triggers:} Identify and limit intake of hot flush triggers such as caffeine, alcohol, tobacco, and spicy foods.
    \item \textbf{Staying active:} Commit to daily physical activity, aiming for at least 30 minutes of moderate exercise.
    \item \textbf{Ensuring calcium-rich nutrition:} Include dairy products, fortified plant milks, leafy greens, and almonds in the daily diet.
    \item \textbf{Practising paced breathing:} Perform slow, deep diaphragmatic breathing (5 to 6 breaths per minute) for 15 minutes twice daily or during the onset of a hot flush.
    \item \textbf{Seeking social support:} Join support groups or speak with friends to share experiences and reduce the feeling of isolation during the menopausal transition.
\end{itemize}

\section*{Sources and Further Reading}
The following reputable organisations provide information on managing menopause without medication:
\begin{itemize}
    \item Australasian Menopause Society (AMS). \textit{Factsheets on non-hormonal and lifestyle management of menopause.} https://www.menopause.org.au/
    \item Jean Hailes for Women's Health. \textit{Resources on menopause symptoms, diet, and lifestyle changes.} https://www.jeanhailes.org.au/
    \item Healthdirect Australia. \textit{Overview of natural and non-medical options for menopause.} https://www.healthdirect.gov.au/
    \item Menopause Practitioner (North American Menopause Society). \textit{Clinical and patient guides on non-hormonal menopause therapies.} https://www.menopause.org/
    \item NHS. \textit{Guide to managing menopause symptoms through lifestyle and self-care.} https://www.nhs.uk/
\end{itemize}
